\subsection{Feature Ablation Study}
\label{sec:ablation}

To understand which CSI feature families actually drive the localization
performance, an ablation was run on the Otaniemi\_small dataset over 22
combinations of the eight candidate feature groups: raw OFDM taps,
time-difference-of-arrival (TDoA), angle-of-arrival (AoA), received
signal strength (RSS), path-loss (PL), per-path delay (Dly), covariance
eigenvalues (CovE), and a LOS-reached indicator (Rch).
For each combination the three regression heads from
Section~\ref{sec:localization} (wKNN, NN, CNN) were re-trained on an
identical train/test split and scored by MAE, RMSE and $P_{90}$.
Table~\ref{tab:ablation} lists the best MAE for each feature combination,
sorted from best to worst.

\begin{table}[htbp]
\caption{Feature-ablation results on Otaniemi\_small.  ``Best MAE'' is
         the minimum across the three regressors; ``Method'' names the
         winner.  Feature families: OFDM = raw OFDM taps,
         TDoA = time-difference of arrival, AoA = angle of arrival,
         RSS = received power, PL = path loss, Dly = per-path delay,
         CovE = covariance eigenvalues, Rch = reached/LOS flag.}
\label{tab:ablation}
\centering
\resizebox{\columnwidth}{!}{%
\begin{tabular}{lrlrl}
\toprule
Combination          & $N_{\mathrm{feat}}$ & Features                      & Best MAE (m) & Method \\
\midrule
Light                & 90                   & TDoA+AoA+RSS+Rch              & \textbf{6.58} & NN Reg \\
No\_OFDM\_Light      & 108                  & TDoA+AoA+RSS+PL+Dly+Rch       & 6.61          & CNN Reg \\
Solo\_AoA            & 36                   & AoA                           & 7.48          & CNN Reg \\
No\_OFDM             & 135                  & TDoA+AoA+RSS+PL+Dly+CovE+Rch  & 7.48          & CNN Reg \\
Geometry             & 81                   & TDoA+AoA+Dly                  & 7.60          & CNN Reg \\
Solo\_RSS            & 9                    & RSS                           & 9.94          & wKNN \\
Radio                & 27                   & RSS+PL+Rch                    & 10.61         & wKNN \\
No\_Reached          & 3681                 & OFDM+TDoA+AoA+RSS+PL+Dly+CovE & 9.47          & NN Reg \\
No\_TDoA             & 3654                 & OFDM+AoA+RSS+PL+Dly+CovE+Rch  & 9.64          & NN Reg \\
No\_Delay            & 3681                 & OFDM+TDoA+AoA+RSS+PL+CovE+Rch & 9.73          & NN Reg \\
No\_CovEig           & 3663                 & OFDM+TDoA+AoA+RSS+PL+Dly+Rch  & 9.83          & NN Reg \\
No\_RSS              & 3681                 & OFDM+TDoA+AoA+PL+Dly+CovE+Rch & 9.94          & NN Reg \\
No\_PathLoss         & 3681                 & OFDM+TDoA+AoA+RSS+Dly+CovE+Rch & 10.02        & NN Reg \\
No\_AoA              & 3654                 & OFDM+TDoA+RSS+PL+Dly+CovE+Rch & 13.34         & NN Reg \\
OFDM\_Radio          & 3573                 & OFDM+RSS+PL                   & 13.83         & NN Reg \\
Solo\_OFDM           & 3555                 & OFDM                          & 13.93         & NN Reg \\
Spectral             & 3555                 & OFDM                          & 14.09         & NN Reg \\
Solo\_CovEig         & 27                   & CovE                          & 32.45         & CNN Reg \\
Solo\_Reached        & 9                    & Rch                           & 50.78         & CNN Reg \\
Solo\_PathLoss       & 9                    & PL                            & 50.87         & CNN Reg \\
Solo\_Delay          & 9                    & Dly                           & 51.19         & CNN Reg \\
Solo\_TDoA           & 36                   & TDoA                          & 52.51         & NN Reg \\
\bottomrule
\end{tabular}%
}
\end{table}

Four observations are worth highlighting:

\paragraph{AoA dominates.}
The 36-feature \texttt{Solo\_AoA} configuration already reaches
$7.48\,\text{m}$ MAE, beating every combination that includes the full
3\,555 raw OFDM taps.  Angles of arrival across the nine BS sites carry
most of the useful geometric signal in a multi-BS outdoor
scene~\cite{3gpp_nr_channel}.

\paragraph{Raw OFDM hurts more than it helps.}
wKNN on \texttt{Solo\_OFDM} reaches only $15.01\,\text{m}$ MAE, roughly
double the compact combinations.  The 3\,555 frequency taps dominate the
Euclidean distance computation used by wKNN and obscure the geometric
features; the NN regressor can partially learn to filter them
(MAE $\sim 10\,\text{m}$) but the full-feature baseline still
\emph{loses} to the 90-feature \texttt{Light} set
($6.58\,\text{m}$) --- a textbook curse-of-dimensionality
effect~\cite{vandermaaten2008tsne}.

\paragraph{Single scalar features are useless alone but useful in combination.}
\texttt{Solo\_TDoA}, \texttt{Solo\_Delay}, \texttt{Solo\_PathLoss} and
\texttt{Solo\_Reached} all collapse to $50$--$56\,\text{m}$ MAE on their
own because each scalar is geometrically ambiguous across the scene.
Yet \emph{removing} any one of them from the full stack measurably
worsens the result (\emph{e.g.}\ removing AoA lifts wKNN from
$7.76$ to $14.82\,\text{m}$), showing that their value lies in the
joint manifold rather than in any individual coordinate.

\paragraph{Sweet spot: the Light configuration.}
The 90-feature \texttt{Light} set (TDoA+AoA+RSS+Rch) attains the best
overall MAE ($6.58\,\text{m}$, NN regression), a 31\,\% improvement on
the 3\,690-feature baseline used in Section~\ref{sec:localization}
($9.53\,\text{m}$), with $97.6\,\%$ fewer features and correspondingly
faster training and inference.
This set is therefore the recommended working point for any follow-on
Otaniemi experiments and is consistent with the feature guidance in the
broader channel-charting literature~\cite{studer2018channel_charting}.
